% Generated draft: requires mathematical and source review.
\documentclass[11pt]{article}
\usepackage[margin=25mm]{geometry}
\usepackage{fontspec}
\setmainfont{Latin Modern Roman}
\usepackage{amsmath,amssymb,mathtools,unicode-math}
\setmathfont{Latin Modern Math}
\usepackage{enumitem,needspace,xurl,hyperref}
\setlength{\parindent}{0pt}
\setlength{\parskip}{5pt}
\setlength{\emergencystretch}{3em}
\newcommand{\N}{\mathbb{N}}
\newcommand{\Z}{\mathbb{Z}}
\newcommand{\Q}{\mathbb{Q}}
\newcommand{\R}{\mathbb{R}}
\newcommand{\C}{\mathbb{C}}
\newcommand{\F}{\mathbb{F}}
\newcommand{\A}{\mathbb{A}}
\newcommand{\PP}{\mathbb{P}}
\newcommand{\E}{\mathbb{E}}
\newcommand{\eps}{\varepsilon}
\newcommand{\dd}{\,\mathrm{d}}
\begin{document}
\section*{UnsolvedMath ID: 30001279}
\textbf{Source problem code:} catalogue entry
\subsection{Definitions and mathematical statement}
A field \(k\) is perfect if every algebraic extension of \(k\) is separable. Let \(X\) be smooth and connected over a perfect field \(k\). The notation \(H_i^{KS}(X,\Z)\) denotes the groups called Kato--Suslin homology in the supplied record. The record does not define these groups or specify the construction and conventions needed to interpret them; consequently, a self-contained mathematical formulation cannot be recovered from the supplied evidence alone.

The proposed assertion, with the source notation retained, is that for every perfect field \(k\) and every such \(X\),
\[
H_i^{KS}(X,\Z)=0\quad\text{for every integer }i>0,
\qquad H_0^{KS}(X,\Z)\cong\Z.
\]
The degree-zero claim is expressed as an isomorphism of abelian groups; no particular canonical identification is supplied. The title refers to varieties, but the statement does not explicitly specify finite type, separatedness, or nonemptiness. These qualifications, together with the definition of \(H_i^{KS}\), require source clarification. The displayed assertion is a statement-only transcription of the proposed conjecture, not an established theorem.
\subsection{Short English statement}
The proposed conjecture says that the integral Kato--Suslin homology of every smooth connected object over a perfect field vanishes in positive degrees and is isomorphic to \(\Z\) in degree zero. The supplied record leaves the homology construction and the precise class of objects unspecified, so the formulation requires mathematical review.
\subsection{Sources}
\begin{itemize}
\item Local UnsolvedMath registry: supplied record with permanent ID \(30001279\), exact source problem code \texttt{catalogue entry}, titled \emph{Kato--Suslin Homology of Smooth Connected Varieties}, in the Oberwolfach Reports Open Problems collection. The proposed year is 2009; no individual proposer is supplied.
\item \emph{Algebraische Zahlentheorie} (2009), Oberwolfach Reports, pp.\ 1667--1730, DOI 10.4171/owr/2009/30. \url{https://doi.org/10.4171/owr/2009/30}. The supplied verification metadata labels the statement \texttt{corrected\_verified} and reports checking it against Conjecture 2.1 of the official report, retaining smoothness, connectedness, positive-degree vanishing, and the degree-zero claim. This verification is reported from the record and has not been independently performed here.
\item Supplied literature reference labeled ``Geisser on Kato's conjecture'': \url{https://doi.org/10.1515/crelle.2010.050}. No more detailed bibliographic information or precise supporting result is supplied.
\item Supplied literature assessment dated exactly \texttt{2026-08-21}: the record labels the problem open and describes it as a generalization of Kato's conjecture for which the extraction gives no complete proof and the available literature does not establish the claimed generality. This is supplied dated curation derived from the OWR record, not an independently verified literature assessment; time-sensitive status and cited claims require independent verification.
\end{itemize}
\end{document}
